\section{The problem selected from the manuscript}
\label{sec:context}
\subsection{Notation and the boundary-value convention}
We use the decreasing-index convention
\begin{equation}\label{eq:F}
 F_{a,b}(z)=\Li_{a,b}(z,1)
 =\sum_{n=1}^{\infty}\frac{H_{n-1}^{(b)}}{n^a}z^n,
 \qquad H_m^{(b)}=\sum_{j=1}^{m}j^{-b},\quad H_0^{(b)}=0.
\end{equation}
Initially $|z|<1$ and $a,b>0$. The positive integral below continues $F$ to $\C\setminus[1,\infty)$. In particular,
\[
 g_{a,b}=\Imag F_{a,b}(\ii)
\]
always denotes this analytic value. The ordinary Fourier series at $\ii$ need not converge when the coefficients in \eqref{eq:F} fail to tend to zero. Nothing in the following proofs requires it to converge.

Set
\begin{equation}\label{eq:Euler}
 A_k(a,b)=\frac{H_{2k}^{(b)}}{(2k+1)^a},\qquad
 d_j(a,b)=\sum_{k=0}^{j}(-1)^k\binom jk A_k(a,b),\qquad
 E_N(a,b)=\sum_{j=0}^{N-1}\frac{d_j(a,b)}{2^{j+1}}.
\end{equation}
Thus $A_0=d_0=E_1=0$. We also use $E_0=0$ and distinguish the unscaled and scaled errors
\begin{equation}\label{eq:errors}
 e_N=E_N-g_{a,b},\qquad R_N=2^N e_N\quad(N\ge0).
\end{equation}
That distinction matters: $e_N$ decreases, but $R_N$ need not do so.

The auxiliary parameter $b$ is a real order, not a depth. The function has depth two throughout. We retain the standard notation
\[
 \beta(s)=\sum_{k\ge0}\frac{(-1)^k}{(2k+1)^s},\qquad
 \eta(s)=\sum_{n\ge1}\frac{(-1)^{n-1}}{n^s}
\]
when the series converge, and use their analytic continuations otherwise. The elementary polylogarithm and gamma identities used here are consistent with the principal conventions in the DLMF \cite{DLMFpoly,DLMFgamma}.

\subsection{What was already proved, and what remained open}
At the pinned snapshot, the manuscript's signed-kernel and real-order chapters establish one-sided Euler convergence, a constant-one estimate for $a\ge1$, and several parameter-dependent bounds for $0<a<1$. The subcritical chapter proves the sharp constant
\begin{equation}\label{eq:critical}
 C_{\rm crit}=\frac\pi4+\frac{\log2}{2}
\end{equation}
uniformly on $a,b>0$, $a+b\le1$. It also proves that
\begin{equation}\label{eq:axisC}
 C(b)=\beta(b)+2^{-b}\eta(b)
\end{equation}
has a unique nondegenerate global maximum, at an order $b_*\in(1,2)$, and that this is the supremum of $-2g_{a,b}$ over all positive orders. Those results are inherited, not discoveries of this continuation \cite{RepoSubcritical,RepoEuler}.

The research programme nevertheless explicitly leaves a universal real-order Euler estimate above the critical triangle open. The warning is correct: bounding the Gaussian value only controls $N=1$, since $2(E_1-g)=-2g$. It does not control later scaled errors without another theorem. This is precisely the gap treated here \cite{RepoDiscovery}.

The new contribution is the uniform estimate for \emph{every} $N\ge2$. It makes the inherited axis maximum decisive for the full Euler problem. The sharp constant is not guessed by maximizing a finite numerical table. An analytic reduction and exact polynomial certificates exclude all later truncations from competing with the first one.

\subsection{Main result and its logical components}
\begin{theorem}[Sharp unrestricted real-order Euler bound]\label{thm:main}
For all $a,b>0$ and integers $N\ge1$,
\begin{equation}\label{eq:sharp-main}
 0<E_N(a,b)-g_{a,b}<\Cs\,2^{-N},
 \qquad \Cs=\max_{b>0}C(b).
\end{equation}
The constant $\Cs$ is the least constant valid uniformly on this domain. More strongly,
\begin{equation}\label{eq:nine-eighths}
 0<E_N(a,b)-g_{a,b}<\frac98\,2^{-N}\qquad(N\ge2).
\end{equation}
The sharp supremum in \eqref{eq:sharp-main} is approached with $N=1$, $b=b_*$ and $a\downarrow0$. If the axis $a=0$, $b>0$ is included, the non-strict sharp estimate has equality only at $(a,b,N)=(0,b_*,1)$.
\end{theorem}
The proof of \eqref{eq:nine-eighths} occupies Sections~\ref{sec:kernel}--\ref{sec:ineq}; Section~\ref{sec:sharp} identifies the optimum. Section~\ref{sec:cert} gives exact rational enclosures, not just numerical approximations. Section~\ref{sec:identities} derives an independent family of all-order identities from the error generating function.

The underlying Euler transform is classical. Positive-measure arguments for sequence acceleration go back well beyond this project; the framework of Cohen--Rodriguez Villegas--Zagier is especially relevant \cite{CRVZ}. Euler-type transformations involving harmonic numbers also have a substantial literature \cite{Boyadzhiev}. The present claim is a specific uniform theorem for the harmonic depth-two family, not the invention of Euler acceleration or of positive-kernel methods.
